\title{Group Sequence Policy Optimization}

\begin{abstract}
We present Group Sequence Policy Optimization (GSPO), a robust and efficient reinforcement learning method tailored for training large language models. In contrast to existing approaches, which typically utilize importance ratios computed at the token level, GSPO establishes importance ratios over entire sequence probabilities and conducts clipping, reward assignment, and optimization at the sequence granularity. Through empirical evaluation, we show that GSPO delivers higher training efficiency and better performance relative to the GRPO algorithm, brings notable stability to Mixture-of-Experts (MoE) reinforcement learning processes, and offers opportunities to streamline RL system architecture. These advantages have significantly contributed to the advances observed in the latest Qwen3 models.
\end{abstract}

\section{Introduction}
\label{sec:intro}

Reinforcement learning (RL) has established itself as a cornerstone for enhancing the scalability of language models \citep{o1,dpsk-r1,qwq32b,qwen3}. Leveraging RL on a large scale enables language models to address complex tasks—including advanced mathematics and programming—by supporting more extended and intricate reasoning sequences.

Effectively scaling RL with increased computational resources fundamentally depends on preserving stable and resilient training dynamics. Nevertheless, leading RL methods such as GRPO \citep{grpo} encounter pronounced instability when applied to massive language models, frequently causing sudden and irreversible model degeneration \citep{qwen3, minimax-m1}. Such instability severely limits the advancement of language model capabilities through ongoing RL-driven improvements.

In this work, we demonstrate that the instability observed in GRPO arises primarily from the flawed application and subsequent invalidation of importance sampling weights within the structure of the algorithm. This misstep introduces substantial variance into the training process, causing noise to build up as response sequences grow longer, an effect which is further exacerbated by the use of the clipping mechanism, ultimately resulting in model collapse.

To overcome these significant challenges, we introduce \textbf{Group Sequence Policy Optimization (GSPO)}, a novel reinforcement learning framework tailored for training large language models. The central contribution of GSPO is a theoretically justified formulation of the importance ratio that leverages sequence likelihoods \citep{click}, thereby restoring adherence to the foundational principle underlying importance sampling. Furthermore, GSPO employs normalized rewards calculated as the advantages of various responses to a particular query, thereby guaranteeing the consistency between reward assignment at the sequence level and the optimization procedure.

Our empirical results indicate that GSPO surpasses GRPO by a considerable margin in terms of training stability, efficiency, and overall performance. Notably, GSPO intrinsically addresses the instability issues that typically arise when training large Mixture-of-Experts (MoE) models in reinforcement learning, thereby removing the necessity for intricate stabilization methods and opening avenues to streamline RL system design. These advantages have played a pivotal role in achieving the marked performance gains seen in the most recent Qwen3 models. Looking ahead, we consider GSPO to be a dependable and scalable foundation for future progress in large-scale RL-based language model training.

\section{Preliminaries}

\paragraph{Notation}
Throughout this work, we represent an autoregressive language model with parameters $\theta$ as the policy $\pi_\theta$. The variable $x$ refers to a single query, while $\mathcal{D}$ stands for the collection of all queries. For a response $y$ corresponding to a particular query $x$, the conditional probability under policy $\pi_\theta$ is written as $\pi_\theta (y | x)=\prod_{t=1}^{|y|} \pi_\theta (y_t | x, y_{<t})$, where $|y|$ indicates the total number of tokens contained in $y$. A verifier $r$ assigns a score to each query-response pair $(x, y)$, with the resulting reward $r(x, y)$ taking values within $[0, 1]$.

\paragraph{Proximal Policy Optimization (PPO)} Using samples generated from the old policy $\pi_{\theta_\text{old}}$, PPO \citep{ppo} constrains the policy update within a proximal region of the old policy through the clipping mechanism. Specifically, PPO employs the following objective for policy optimization (we omit the KL regularization term hereinafter for brevity, as it is not the focus of this paper):

\begin{align}
\mathcal{J}_\text{PPO}(\theta) = \mathbb{E}_{ x \sim \mathcal{D},\, y \sim \pi_{\theta_\text{old}}( \cdot | x) }
\left[ \frac{1}{|y|} \sum_{t=1}^{|y|} 
\min \left( w_{t}(\theta) \widehat{A}_{t},  \, \mathrm{clip} \left( w_{t}(\theta), 1 - {\varepsilon}, 1 + {\varepsilon}\right) \widehat{A}_{t} \right)
\right],
\end{align}

where the token importance ratio $w_{t}(\theta)$ is given by $
w_{t}(\theta) = \frac{ \pi_{\theta} (y_{t} | x, y_{<t}) }{ \pi_{\theta_\text{old}} (y_{t} | x,y_{<t})} $, $\widehat{A}_{t}$ refers to the estimated advantage of $y_t$ provided by a separate value model, and $\varepsilon$ denotes the clipping threshold applied to the importance ratios.

A principal obstacle in applying PPO in real-world scenarios is its significant dependence on the value model. Typically, the value model is nearly equal in size to the policy model, which leads to increased computational and memory demands. Moreover, the overall success of the algorithm is tightly linked to the accuracy of value estimation. In practice, obtaining dependable value estimates is intrinsically difficult, and this challenge is further amplified when scaling to generate longer outputs or to tackle more intricate tasks.

\paragraph{Group Relative Policy Optimization (GRPO)} GRPO \citep{grpo} bypasses the need for the value model by computing the relative advantage of each response within a group of responses to the same query. Specifically, GRPO optimizes the following objective:

\begin{align}
\label{equ:grpo}
\mathcal{J}_\text{GRPO}(\theta) = \mathbb{E}_{ x \sim \mathcal{D},\, \{y_i\}_{i=1}^G \sim \pi_{\theta_\text{old}}( \cdot | x) }
\left[ \frac{1}{G} \sum_{i=1}^{G} \frac{1}{|y_i|} \sum_{t=1}^{|y_i|} 
\min \left( w_{i,t}(\theta) \widehat{A}_{i,t},  \, \mathrm{clip} \left( w_{i,t}(\theta), 1 - {\varepsilon}, 1 + {\varepsilon}\right) \widehat{A}_{i,t} \right)
\right],
\end{align}

where $G$ is the number of generated responses to each query $x$ (i.e., the group size), and the importance ratio $w_{i,t}(\theta)$ and advantage $\widehat{A}_{i,t}$ of token $y_{i,t}$ are:

\begin{align}
    w_{i,t}(\theta)=\frac{ \pi_{\theta} (y_{i,t} | x, y_{i,<t}) }{ \pi_{\theta_\text{old}} (y_{i,t} | x,y_{i,<t})},\quad\ 
    \widehat{A}_{i,t} = \widehat{A}_{i} = \frac{r(x, y_i) - \mathrm{mean} \left( \{ r(x, y_i) \}_{i=1}^G \right) }{ \mathrm{std} \left( \{ r(x, y_i) \}_{i=1}^G \right) },
\end{align}

respectively, where all the tokens in $y_i$ share the same advantage as $\widehat{A}_{i}$.

\section{Motivation}
\label{sec:motivation}

As models increase in size, adopt more sparse architectures (such as Mixture-of-Experts), and generate longer responses, employing larger rollout batch sizes becomes critical for efficient use of computational resources during reinforcement learning. To enhance sample efficiency, practitioners typically divide these extensive rollout batches into several mini-batches when performing gradient computations. This approach inherently leads to an off-policy learning context, as sampled responses $y$ originate from a previous policy $\pi_{\theta_\text{old}}$ instead of the current policy $\pi_\theta$ under optimization. Such off-policy scenarios necessitate mechanisms like the clipping strategies found in PPO and GRPO, which serve to limit the influence of samples that significantly diverge from the present policy during gradient calculations.

While mechanisms like clipping aim to manage this off-policy discrepancy, we identify a more fundamental issue in GRPO: \textit{its objective is ill-posed}. This problem becomes particularly acute when training large models on long-response tasks, leading to catastrophic model collapse. The ill-posed nature of the GRPO objective stems from a misapplication of importance sampling weights. The principle of importance sampling is to estimate the expectation of a function $f$ under a target distribution $\pi_\text{tar}$ by re-weighting samples drawn from a behavior distribution $\pi_\text{beh}$:

\begin{align}
\mathbb{E}_{z \sim \pi_\text{tar}} \left[ f(z) \right] = \mathbb{E}_{z \sim \pi_\text{beh}} \left[ \frac{ \pi_\text{tar} (z) }{ \pi_\text{beh} (z)} \, f(z) \right].
\end{align}

A central aspect of this approach is its dependence on taking an average over a large number of samples ($N \gg 1$) drawn from the behavioral policy $\pi_\text{beh}$, enabling the importance weighting term $\frac{ \pi_\text{tar} (z) }{ \pi_\text{beh} (z)}$ to reliably adjust for differences between the behavior and target distributions.

On the other hand, GRPO calculates the importance weight $\frac{ \pi_{\theta} (y_{i,t} | x, y_{i,<t}) }{ \pi_{\theta_\text{old}} (y_{i,t} | x,y_{i,<t})}$ at every token timestep $t$. Because this is computed from only one sampled token $y_{i,t}$ per conditional next-token distribution $\pi_{\theta_\text{old}}( \cdot | x, y_{i, <t})$, the importance weight does not effectively fulfill its corrective function. This procedure instead injects substantial high-variance noise into the gradient estimates during training, with this variability compounding over extended sequences and further intensified by the implementation of the clipping operation. In practice, we have observed this frequently culminates in catastrophic model collapse, from which recovery is extremely rare. Attempts to resume training after collapse—including rolling back to earlier model checkpoints, extensive hyperparameter searches (such as adjusting clipping thresholds), increasing the maximum generation length, or changing the RL query process—have proved ineffective.

These findings underscore a more profound flaw in the architecture of GRPO. The ineffectiveness of token-level importance weighting underlines a central insight: \textbf{the granularity of the optimization target should correspond to that of the reward signal}. Since the reward function evaluates the full sequence as a unit, implementing off-policy corrections at the token scale appears inadequate. This realization leads us to abandon token-level objectives in favor of approaches that employ importance weighting and carry out optimization over complete sequences.

\section{Algorithm}

\subsection{GSPO: Group Sequence Policy Optimization}

Although GRPO faces issues with the token-level importance weight $\frac{ \pi_{\theta} (y_{i,t} | x, y_{i,<t}) }{ \pi_{\theta_\text{old}} (y_{i,t} | x,y_{i,<t})}$, it is noteworthy that, in language generation tasks, the \textit{sequence-level} importance weight $\frac{ \pi_{\theta} (y | x) }{ \pi_{\theta_\text{old}} (y | x)}$ possesses an explicit theoretical interpretation. Specifically, it quantifies the degree to which a response $y$, drawn according to $\pi_{\theta_\text{old}} (\cdot | x)$, diverges from the distribution $\pi_{\theta} (\cdot | x)$. This correspondence is naturally coherent with the sequence-level reward framework and additionally provides a valuable signal for interpreting the role of the clipping mechanism.

Based on this straightforward observation, we propose the \textbf{Group Sequence Policy Optimization (GSPO)} algorithm. GSPO employs the following sequence-level optimization objective:

\begin{align}
\mathcal{J}_\text{GSPO} (\theta) =
\mathbb{E}_{ x \sim \mathcal{D},\, \{y_i\}_{i=1}^G \sim \pi_{\theta_\text{old}}( \cdot | x) }
\left[ 
\frac{1}{G} \sum_{i=1}^{G}
\min \left( s_{i}(\theta)  \widehat{A}_{i},  \, \mathrm{clip} \left( s_{i}(\theta), 1 - {\varepsilon}, 1 + {\varepsilon} \right) \widehat{A}_{i} \right) 
\right],
\label{equ:gspo}
\end{align}

where we adopt the group-based advantage estimation:

\begin{align}
\widehat{A}_{i} = \frac{r(x, y_i) - \mathrm{mean} \left( \{ r(x, y_i) \}_{i=1}^G \right) }{ \mathrm{std} \left( \{ r(x, y_i) \}_{i=1}^G \right) },
\end{align}

and define the importance ratio $s_{i}(\theta)$ based on sequence likelihood \citep{click}:

\begin{align}
s_{i}(\theta) = \left( \frac{ \pi_{\theta} (y_i | x) }{ \pi_{\theta_\text{old}} (y_i | x)} \right)^{\frac{1}{|y_i|}}
=
\exp \left( \frac{1}{|y_i|} \sum_{t=1}^{|y_i|} \log \frac{ \pi_{\theta} (y_{i,t} | x, y_{i,<t}) }{ \pi_{\theta_\text{old}} (y_{i,t} | x,y_{i,<t})} \right).
\end{align}

Accordingly, GSPO implements clipping at the response level rather than per token, effectively filtering out excessively ``off-policy'' samples from the gradient calculations, which aligns with both the reward evaluation and optimization processes at the sequence scale. Additionally, we use length normalization for $s_{i}(\theta)$ to mitigate variance and maintain $s_{i}(\theta)$ within a consistent numeric interval. Without this normalization, substantial likelihood changes in just a few tokens could cause pronounced shifts in the overall sequence-level importance ratio, necessitating different clipping thresholds depending on response length. It is also important to highlight that the scale of clipping ranges in GSPO contrasts—often by orders of magnitude—with prior methods such as GRPO, due to the fundamentally different definitions used for importance ratios.

\subsection{Gradient Analysis}
\label{subsec:gradient}

We can derive the gradient of the GSPO objective as follows (clipping is omitted for brevity):

\begin{align}
\nabla_{\theta} \mathcal{J}_\text{GSPO} (\theta)
=&\ 
\nabla_{\theta} \mathbb{E}_{ x \sim \mathcal{D},\, \{y_i\}_{i=1}^G \sim \pi_{\theta_\text{old}}( \cdot | x) }
\left[ 
\frac{1}{G} \sum_{i=1}^{G} s_{i}(\theta) \widehat{A}_{i}
\right] \\
= &\ 
\mathbb{E}_{ x \sim \mathcal{D},\, \{y_i\}_{i=1}^G \sim \pi_{\theta_\text{old}}( \cdot | x) }
\left[ 
\frac{1}{G} \sum_{i=1}^{G}
s_{i}(\theta) \widehat{A}_{i}
\cdot \nabla_{\theta} \log s_{i}(\theta)
\right] \\
=&\ 
\mathbb{E}_{ x \sim \mathcal{D},\, \{y_i\}_{i=1}^G \sim \pi_{\theta_\text{old}}( \cdot | x) }
\left[ 
\frac{1}{G} \sum_{i=1}^{G} 
\left( \frac{ \pi_{\theta} (y_i | x) }{ \pi_{\theta_\text{old}} (y_i | x)} \right)^{\frac{1}{|y_i|}} \widehat{A}_{i} 
\cdot \frac{1}{|y_i|} \sum_{t=1}^{|y_i|} \nabla_{\theta} \log \pi_{\theta} (y_{i,t} | x, y_{i,<t}) 
\right].
\label{equ:gspo_grad}
\end{align}

For comparison, the gradient of the GRPO objective is as follows (note that $\widehat{A}_{i,t} = \widehat{A}_{i}$):

\begin{align}
\nabla_{\theta} \mathcal{J}_\text{GRPO} (\theta) 
=&\ 
\nabla_{\theta} \mathbb{E}_{ x \sim \mathcal{D},\, \{y_i\}_{i=1}^G \sim \pi_{\theta_\text{old}}( \cdot | x) }
\left[ \frac{1}{G} \sum_{i=1}^{G} \frac{1}{|y_i|} \sum_{t=1}^{|y_i|} 
w_{i,t}(\theta) \widehat{A}_{i,t}
\right] \\
=&\
\mathbb{E}_{ x \sim \mathcal{D},\, \{y_i\}_{i=1}^G \sim \pi_{\theta_\text{old}}( \cdot | x) }
\left[ \frac{1}{G} \sum_{i=1}^{G} \widehat{A}_{i} 
\cdot \frac{1}{|y_i|} \sum_{t=1}^{|y_i|} 
\frac{ \pi_{\theta} (y_{i,t} | x, y_{i,<t}) }{ \pi_{\theta_\text{old}} (y_{i,t} | x,y_{i,<t})} 
\nabla_{\theta} \log \pi_{\theta} (y_{i,t} | x, y_{i,<t})  
\right].
\end{align}

Consequently, the primary difference between GSPO and GRPO is rooted in \textit{the method by which each approach applies weighting to the gradients of token log likelihoods}. In the case of GRPO, each token is assigned a weight based on its corresponding ``importance weight'' $\frac{ \pi_{\theta} (y_{i,t} | x, y_{i,<t}) }{ \pi_{\theta_\text{old}} (y_{i,t} | x,y_{i,<t})}$. These non-uniform weights are far from trivial, as they can fall within the range $(0, 1+\varepsilon]$ when $\widehat{A}_{i} > 0$ or $[1-\varepsilon, +\infty)$ when $\widehat{A}_{i} < 0$, and their heterogeneity may accumulate throughout training, potentially giving rise to unstable or unforeseen effects. On the other hand, GSPO removes this source of instability by treating all tokens in a response identically, assigning them uniform weights.

\subsection{GSPO-token: A Token-level Objective Variant}

In scenarios like multi-turn RL, we may desire a finer-grained advantage adjustment than the sequence level. To this end, we introduce a token-level objective variant of GSPO, namely \textbf{GSPO-token}, to allow token-wise advantage customization:

\begin{align}
\mathcal{J}_\text{GSPO-token}(\theta)
=
\mathbb{E}_{ x \sim \mathcal{D},\, \{y_i\}_{i=1}^G \sim \pi_{\theta_\text{old}}( \cdot | x) }
\left[ 
\frac{1}{G} \sum_{i=1}^{G} \frac{1}{|y_i|} \sum_{t=1}^{|y_i|} 
\min \left( s_{i,t}(\theta) \widehat{A}_{i,t},  \, \mathrm{clip} \left( s_{i,t}(\theta), 1 - {\varepsilon}, 1 + {\varepsilon} \right) \widehat{A}_{i,t} \right) 
\right],
\label{equ:gspo-token}
\end{align}

where

\begin{align}
s_{i,t}(\theta) 
= \mathrm{sg} \left[ s_{i}(\theta) \right]  \cdot \frac{ \pi_{\theta} (y_{i,t} | x, y_{i,<t}) }{ \mathrm{sg} \left[ \pi_{\theta} (y_{i,t} | x, y_{i,<t}) \right] },
\end{align}

and $\mathrm{sg}[\cdot]$ denotes only taking the numerical value but stopping the gradient, corresponding to the \texttt{detach} operation in PyTorch. The gradient of GSPO-token can be derived as:

\begin{align}
\nabla_{\theta} \mathcal{J}_\text{GSPO-token}(\theta)
=&\ 
\nabla_{\theta} \mathbb{E}_{ x \sim \mathcal{D},\, \{y_i\}_{i=1}^G \sim \pi_{\theta_\text{old}}( \cdot | x) }
\left[ 
\frac{1}{G} \sum_{i=1}^{G}
\frac{1}{|y_i|} \sum_{t=1}^{|y_i|} 
s_{i,t}(\theta) \widehat{A}_{i,t}
\right] \\
=&\ 
\mathbb{E}_{ x \sim \mathcal{D},\, \{y_i\}_{i=1}^G \sim \pi_{\theta_\text{old}}( \cdot | x) }
\left[ 
\frac{1}{G} \sum_{i=1}^{G} s_i (\theta) \cdot \frac{1}{|y_i|} \sum_{t=1}^{|y_i|} 
\widehat{A}_{i,t} \frac{ \nabla_{\theta} \pi_{\theta} (y_{i,t} | x, y_{i,<t}) }{ \pi_{\theta} (y_{i,t} | x, y_{i,<t}) }
\right] \\
=&\ 
\mathbb{E}_{ x \sim \mathcal{D},\, \{y_i\}_{i=1}^G \sim \pi_{\theta_\text{old}}( \cdot | x) }
\left[ 
\frac{1}{G} \sum_{i=1}^{G}  \left( \frac{ \pi_{\theta} (y_i | x) }{ \pi_{\theta_\text{old}} (y_i | x)} \right)^{\frac{1}{|y_i|}}
\cdot \frac{1}{|y_i|} \sum_{t=1}^{|y_i|}
\widehat{A}_{i,t} \nabla_{\theta} \log \pi_{\theta} (y_{i,t} | x, y_{i,<t})  
\right].
\label{equ:gspo-token_grad}
\end{align}

Observe that the factor $\frac{ \pi_{\theta} (y_{i,t} | x, y_{i,<t}) }{ \mathrm{sg} \left[ \pi_{\theta} (y_{i,t} | x, y_{i,<t}) \right] }$ evaluates to 1, leading to $s_{i,t}(\theta)$ being numerically identical to $s_{i}(\theta)$. By juxtaposing Equation~\eqref{equ:gspo} with~\eqref{equ:gspo-token} and Equation~\eqref{equ:gspo_grad} with~\eqref{equ:gspo-token_grad}, we see that both GSPO-token and GSPO yield the same results in terms of the optimization objective, the clipping condition, and the theoretical gradient under the assumption that all tokens in the sequence $y_i$ share the same advantage value (that is, $\widehat{A}_{i,t} = \widehat{A}_{i}$). Nevertheless, GSPO-token offers the added capability of assigning distinct advantages at the token level, enabling greater flexibility.

\section{Experiments and Discussion}

\subsection{Empirical Results}

We utilize a cold-start configuration based on Qwen3-30B-A3B-Base for fine-tuning and provide both training reward trajectories and performance metrics on three benchmarks: AIME'24 (measured by average Pass@1 over 32 samples), LiveCodeBench (from 202410 to 202502, averaged Pass@1 over 8 samples), and CodeForces (assessed via Elo Rating). For reinforcement learning training, each rollout batch is subdivided into four mini-batches to facilitate gradient updates. Within GSPO, the left and right clipping boundaries in Equation~\eqref{equ:gspo} are assigned values of 3e-4 and 4e-4, respectively. As a baseline, we employ GRPO, tuning its left and right clipping parameters in Equation~\eqref{equ:grpo} to 0.2 and 0.27, respectively, to guarantee a robust comparison. It should be emphasized that GRPO requires the use of the Routing Replay approach to ensure proper convergence during MoE RL—a method examined further in \S~\ref{subsec:routing_replay}—while \textit{GSPO eliminates the reliance on this training strategy}.

As illustrated in Figure~\ref{fig:results}, the training process using GSPO exhibits consistent stability. Notably, \textbf{GSPO enables ongoing enhancements in performance by scaling up training compute, periodically refreshing the query set, and increasing the generation length}. Additionally, GSPO consistently outperforms GRPO in terms of training efficiency, reaching higher training accuracy and superior benchmark results given equivalent training compute and query usage. Lastly, the successful deployment of GSPO in the RL training of the most recent Qwen3 models further substantiates GSPO’s effectiveness in harnessing RL scaling for large language model development.

\subsection{Curious Observation on Clipping Fractions}

One notable difference between GSPO and GRPO is that GSPO applies clipping at the response level, in contrast to GRPO's token-wise clipping. As illustrated in Figure~\ref{fig:clipping}, GSPO clips a far greater proportion of tokens—by about two orders of magnitude—relative to GRPO (and this substantial difference persists even when the clipping ranges are modified). Nevertheless, even though GSPO excludes many more tokens from training or gradient estimation due to its response-level clipping, it nonetheless demonstrates greater training efficiency than GRPO. This seemingly paradoxical result—where discarding a larger portion of tokens leads to improved training outcomes—suggests that GRPO's gradient calculations at the token level introduce considerable noise, which reduces sample efficiency. Conversely, GSPO's sequence-level strategy delivers a more robust and informative learning signal.

\subsection{Benefit of GSPO for MoE Training}
\label{subsec:routing_replay}

\paragraph{Background}
Relative to the RL training process applied to dense models, the inherently sparse activation in MoE architectures introduces distinct challenges to training stability. Notably, employing the GRPO algorithm reveals that \textit{expert-activation volatility} in MoE models can significantly disrupt the RL training's ability to reach convergence. Specifically, we observe that a few gradient steps can lead to dramatic shifts in the set of experts engaged for identical responses. Taking the 48-layer Qwen3-30B-A3B-Base model as an illustration, approximately 10\% of the experts selected by the updated policy $\pi_{\theta}$ differ from those chosen by the preceding policy $\pi_{\theta_\text{old}}$ for the same sample after each RL gradient step. Such instability intensifies with deeper MoE models, causing the token-level importance weights $w_{i,t}(\theta) = \frac{ \pi_{\theta} (y_{i,t} | x, y_{i,<t}) }{ \pi_{\theta_\text{old}} (y_{i,t} | x,y_{i,<t})}$ to become highly erratic and unreliable, as elaborated in \S~\ref{sec:motivation} and~\ref{subsec:gradient}. This dynamic impedes the regular convergence dynamics of RL training.

\paragraph{Our Previous Approach} In our earlier work, we addressed this problem using the \textbf{Routing Replay} training method. In detail, we store the set of experts selected by $\pi_{\theta_\text{old}}$ and subsequently use these identical routing decisions in $\pi_{\theta}$ while evaluating the importance weights $w_{i,t}(\theta) = \frac{ \pi_{\theta} (y_{i,t} | x, y_{i,<t}) }{ \pi_{\theta_\text{old}} (y_{i,t} | x,y_{i,<t})}$. By doing so, both $\pi_{\theta} (y_{i,t} | x, y_{i,<t})$ and $\pi_{\theta_\text{old}} (y_{i,t} | x,y_{i,<t})$ utilize an identical subset of the network for any given token $y_{i,t}$, leading to improved consistency in token-wise importance ratios. This practice ensures that training consistently updates the same network paths over gradient steps. As illustrated in Figure~\ref{fig:routing_replay}, Routing Replay is a critical component in facilitating the stable convergence of GRPO training for MoE architectures.

\paragraph{Benefit of GSPO} While Routing Replay allows for successful GRPO training of MoE models, it does so at the expense of increased memory usage, higher communication costs, and a potential reduction in the model's effective capacity due to the repeated use of routing configurations. GSPO, in contrast, as depicted in Figure~\ref{fig:results}, dispenses with the reliance on Routing Replay, enabling conventional computation of importance ratios $s_i(\theta)$. This results in stable optimization and reliable convergence. The central idea underlying GSPO is its exclusive dependence on the sequence-level likelihood (i.e., $\pi_{\theta} (y_i | x)$), as opposed to the likelihoods of individual tokens (i.e., $\pi_{\theta} (y_{i,t} | x, y_{i,<t})$), making it robust to token-wise variability. Because the MoE architecture inherently preserves its language modeling ability, the sequence likelihood remains relatively stable. Thus, GSPO inherently mitigates the issue of expert-activation instability prevalent in MoE models, eliminating the necessity for elaborate mechanisms like Routing Replay. As a result, this method streamlines and stabilizes training and empowers the model to exploit its entire capacity, unimpeded by arbitrary restrictions.

\subsection{Benefit of GSPO for RL Infrastructure}

Due to the differences in numerical precision between training engines such as Megatron and inference engines like SGLang or vLLM, it is common practice to recompute the likelihoods of generated responses under the previous policy $\pi_{\theta_\text{old}}$ using the training engine. In contrast, GSPO relies on sequence-level likelihoods for optimization, which are generally less sensitive to such precision inconsistencies than token-level likelihoods. As a result, GSPO enables the direct use of the likelihood outputs from the inference engine during optimization, eliminating the necessity to recalculate them with the training engine. This feature is particularly advantageous for applications such as partial rollout, multi-turn RL, and systems that decouple training from inference.

\section{Conclusion}

We introduce Group Sequence Policy Optimization (GSPO), a novel reinforcement learning technique tailored for large language model training. GSPO adheres to the fundamental concept of importance sampling by establishing importance ratios via sequence likelihood and subsequently employs sequence-wise clipping, rewarding, and optimization operations. When compared to GRPO, GSPO achieves significantly greater stability, training efficiency, and overall performance. It is especially effective for large-scale reinforcement learning in Mixture-of-Experts (MoE) models, and serves as the underpinning for the remarkable advancements seen in the most recent Qwen3 models. As GSPO offers a robust, scalable framework, we intend to further expand reinforcement learning research utilizing this method, anticipating substantial progress in the development of artificial intelligence.

\end{document}